% Paper 1 (Petty's projection-volume conjecture): abstract and introduction lines 1-47
\begin{abstract}
We prove that ellipsoids uniquely minimize the volume of the projection body
among convex bodies of fixed volume in every dimension at least four.
This proves Petty's projection-volume conjecture in these dimensions.
\end{abstract}
\section{Introduction}
\label{sec:introduction}

For a convex body $K\subset\R^n$, its projection body $\Pi K$ is the
origin-symmetric convex body whose support function is
\[
 h_{\Pi K}(u)=\operatorname{vol}_{n-1}
       \bigl(\operatorname{proj}_{u^\perp}K\bigr),
       \qquad u\in S^{n-1}.
\]
Here a convex body is compact, convex, and has nonempty interior, and
$h_L(x)=\max_{y\in L}\ip{x}{y}$ denotes its support function.
Thus $\Pi K$ records the volumes of the orthogonal shadows of $K$ in
all directions. Write $|L|=\operatorname{vol}_n(L)$ and let $B_2^m$
be the Euclidean unit ball in $\R^m$, with volume $\kappa_m$.
The normalized projection volume is
\[
 R_n(K)=\frac{|\Pi K|}{|K|^{n-1}}.
\]
It is invariant under translations and invertible linear maps: the
classical contravariance identity is
$\Pi(TK)=|\det T|T^{-t}\Pi K$; see
\cite[Equation~(1)]{Ludwig2002}. We verify this identity directly in
Section~\ref{sec:conclusion}.

Petty's projection-volume conjecture asserts, for $n\ge3$, that
$R_n$ is minimized precisely by ellipsoids \cite{Petty1971}.
The ball satisfies $\Pi B_2^n=\kappa_{n-1}B_2^n$, so the proposed
minimum is fixed by the Euclidean case. This is an affine isoperimetric
problem in which directional projection data are assembled into a
second body before its volume is measured. The classical Petty
inequality for the \emph{polar} of the projection body concerns a
different functional; see \cite[Section~18.6]{Gardner2002}.
We prove the unpolarized conjecture in every dimension at least four.

\begin{theorem}\label{thm:main}
For every integer $n\ge4$ and every convex body $K\subset\R^n$,
\begin{equation}\label{eq:main}
 \frac{|\Pi K|}{|K|^{n-1}}
 \ge \kappa_{n-1}^{\,n}\kappa_n^{\,2-n}.
\end{equation}
Equality holds if and only if $K=a+TB_2^n$ for some $a\in\R^n$ and some
invertible linear map $T$.
\end{theorem}

The theorem includes bodies without symmetry or boundary regularity.
Its equality assertion is global and imposes no proximity to an

% Paper 2 (product counterexample): abstract, introduction (main.tex lines 38-100) and Corollary (lines 421-431)
\begin{abstract}
The product of two ten-dimensional simplices has larger normalized
projection-body volume than a twenty-dimensional simplex. This gives
a counterexample to Brannen's proposed simplex maximum.
\end{abstract}
\section{Introduction}

Throughout, a convex body in \(\R^d\) is compact and convex with nonempty
interior, and \(d\ge2\). Write \(|K|\) for its \(d\)-dimensional
volume and \(h_C(x)=\max_{y\in C}\langle x,y\rangle\) for the support
function of a compact convex set \(C\). The projection body
\(\Pi_d K\), also written \(\Pi K\) when the dimension is evident, is
defined by
\[
 h_{\Pi K}(u)=\vol_{d-1}\bigl(\proj_{u^\perp}K\bigr),
 \qquad u\in S^{d-1},
\]
extended homogeneously to all of \(\R^d\). Here \(S^{d-1}\) is the
Euclidean unit sphere and \(\proj_{u^\perp}\) is orthogonal projection.
We write \(B_2^j\) for the Euclidean unit ball in \(\R^j\) and
\(\kappa_j=|B_2^j|\).
The functional under consideration is
\begin{equation}\label{eq:functional}
 R_d(K)=\frac{|\Pi K|}{|K|^{d-1}},
 \qquad c_d=\frac{(d+1)d^d}{d!}.
\end{equation}

The extremal values of this affine-invariant functional are classical
questions in convex geometry. Petty established the affine covariance of
the projection-body operator~\cite{Petty}; see also
\cite[Equation~(1)]{Ludwig}. Brannen conjectured in 1996 that simplices
maximize \(R_d\) among all convex bodies~\cite{Brannen}.
The simplex value is \(c_d\), as we verify below. Thus the proposed
simplex upper bound is
\begin{equation}\label{eq:proposed-bound}
 |\Pi K|\le c_d|K|^{d-1}.
\end{equation}
It is the volume of \(\Pi K\) itself that occurs here, not the volume of
its polar: the classical projection inequalities involving the polar
concern a different functional. For \(d\ge3\), Petty's separate minimum
problem asks whether ellipsoids are the only minimizers of the same
functional \(R_d\); see \cite[Section~1]{Saroglou} for these distinctions.

In dimension three, Chen, Feng, Li, Xi, and Xu settle the minimum problem
and also prove
the bound \(R_3(K)\le18=c_3\) for all
three-dimensional convex bodies, with tetrahedra attaining
equality~\cite[Theorems~1.1 and~1.2]{ChenEtAl}. Feng, Hu, Liu, and Xu give
counterexamples to \eqref{eq:proposed-bound} in every dimension
\(d\ge9\), using polytopes with at most \(d+2\)
facets~\cite[Theorem~1.3]{FengEtAl}. Their result already gives a
negative answer to the unrestricted-dimensional simplex-maximum
question. Their construction uses perturbed projections of a cube and
Minkowski's existence theorem~\cite[Section~4]{FengEtAl}. The present
paper supplies a direct Cartesian-product construction and its exact
value: two ten-dimensional simplices suffice.

\begin{theorem}\label{thm:counterexample}
Let \(T_{10}=\conv(0,e_1,\ldots,e_{10})\subset\R^{10}\), where the
\(e_i\) are the coordinate unit vectors, and let
\(K=T_{10}\times T_{10}\subset\R^{20}\). Then
\[
 \frac{R_{20}(K)}{c_{20}}
 =\frac{121\binom{20}{10}}{21\cdot2^{20}}
 =\frac{22\,355\,476}{22\,020\,096}>1.
\]
Thus \(K\) violates \eqref{eq:proposed-bound} in dimension twenty.
\end{theorem}
Iterating the product test gives the following quantitative consequence.

\begin{corollary}\label{cor:exponential-excess}
There exist a constant \(\lambda>1\) and an integer \(N\ge20\) such that, for every
integer \(n\ge N\), some full-dimensional convex polytope
\(K_n\subset\R^n\) satisfies
\[
 \frac{R_n(K_n)}{R_n(T_n)}\ge\lambda^n.
\]
The polytopes \(K_n\) may be chosen as Cartesian products of simplices.
\end{corollary}
